% !TeX program = xelatex
\documentclass[12pt,spanish,letterpaper,oneside]{article}
\PassOptionsToPackage{es-noshorthands}{babel}
\def\documenttitle{Simulación y juego de roles en el análisis del conflicto socioambiental por el trasvase de agua El Cuchillo --- Distrito de Riego 026}
\def\documentsubtitle{Evidencia 3 --- Reporte de simulación de conflicto agrario y negociación deliberativa (Fase 3)}
\def\documentsubject{Responsabilidad Social y Desarrollo Sustentable --- Ingeniero Agrónomo}
\def\documentauthor{Carlos Mauricio Carvajal Coronado}
\def\coursename{Responsabilidad Social y Desarrollo Sustentable}
\def\coursecode{REF\#261013755032}
\def\universityname{Universidad Autónoma de Nuevo León}
\def\universityfaculty{Facultad de Agronomía}
\def\universitydepartment{Ingeniero Agrónomo}
\def\universitydepartmentimage{departamentos/logos-institucionales.pdf}
\def\universitydepartmentimagecfg{height=2.1cm}
\def\universitylocation{General Escobedo / Marín, Nuevo León, México}
\def\authortable{%
	\begin{Form}
	\begin{tabular}{@{}ll@{}}
		\textbf{Estudiante:} & Carlos Mauricio Carvajal Coronado\\[0.35cm]
		\textbf{Matrícula:} & 1977875\\[0.35cm]
		\textbf{Docente Titular:} & Cavazos Galindo Jaime Manuel\\[0.35cm]
		\textbf{Unidad de Aprendizaje:} & Responsabilidad Social y Desarrollo Sustentable (Evidencia 3 \textbar\ 12 Pts)\\[0.5cm]
		\multicolumn{2}{l}{Fecha de realización: 14 de octubre de 2026}\\
		\multicolumn{2}{l}{\universitylocation}
	\end{tabular}
	\end{Form}
}

\usepackage{fontspec}
\setmainfont{Arial}
\setsansfont{Arial}
\AddToHook{file/config.tex/after}[aulatex-rsds]{%
  \def\documentfontsize{12}
  \def\documentinterline{1}
  \def\documentparindent{0}
  \def\documentparskip{5}
  \def\fontdocument{custom}
  \def\fonttypewriter{custom}
  \def\stylecitereferences{custom}
  \def\portraitstyle{style1}
  \def\hfstyle{style1}
  \def\hfwidthwrap{true}
  \def\hfwidthtitle{0.42}
  \def\hfwidthcourse{0.54}
  \def\predocpageromannumber{false}
  \def\pagemarginleft{1.8}
  \def\pagemarginright{1.8}
  \def\pagemargintop{2.7}
  \def\pagemarginbottom{2.5}
  \def\firstpagemargintop{3.8}
  \def\sectionfontsize{\large}
  \def\sectionspacingtop{8}
  \def\sectionspacingbottom{4}
  \def\ssectionfontsize{\normalsize}
  \def\ssectionspacingtop{6}
  \def\ssectionspacingbottom{3}
  \def\defaultimagefolder{UANL/ingeniero-agronomo/responsabilidad-social-y-desarrollo-sustentable/img/}
  \def\usepdfmetadata{false}
}
% Template:     Informe LaTeX
% Documento:    Núcleo del template
% Versión:      DEV
% Codificación: UTF-8
%
% Autor: Pablo Pizarro R.
%        pablo@ppizarror.com
%
% Manual template: [https://latex.ppizarror.com/informe]
% Licencia MIT:    [https://opensource.org/licenses/MIT]

% -----------------------------------------------------------------------------
% CONFIGURACIONES
% -----------------------------------------------------------------------------
% Definiciones previas
\input{src/defs}

% Ajustes usuario
\input{src/config}

% -----------------------------------------------------------------------------
% IMPORTACIÓN DE LIBRERÍAS
% -----------------------------------------------------------------------------
\input{src/env/imports}

% -----------------------------------------------------------------------------
% IMPORTACIÓN DE FUNCIONES Y ENTORNOS
% -----------------------------------------------------------------------------
\input{src/cmd/core}
\input{src/cmd/math}
\input{src/cmd/equation}
\input{src/cmd/image}
\input{src/cmd/title}
\input{src/cmd/other}
\input{src/cmd/column}
\input{src/env/environments}

% -----------------------------------------------------------------------------
% IMPORTACIÓN DE ESTILOS
% -----------------------------------------------------------------------------
\input{src/style/code}
\input{src/style/other}

% -----------------------------------------------------------------------------
% CONFIGURACIÓN INICIAL DEL DOCUMENTO
% -----------------------------------------------------------------------------
\input{src/cfg/init}

% -----------------------------------------------------------------------------
% PORTADA
% -----------------------------------------------------------------------------
\input{src/page/portrait}
\input{src/page/portrait_config}

% -----------------------------------------------------------------------------
% CONFIGURACIÓN DE PÁGINA Y ENCABEZADOS
% -----------------------------------------------------------------------------
\input{src/cfg/page}

% -----------------------------------------------------------------------------
% TABLA DE CONTENIDOS - ÍNDICE
% -----------------------------------------------------------------------------
\input{src/page/index}

% -----------------------------------------------------------------------------
% CONFIGURACIONES FINALES
% -----------------------------------------------------------------------------
\input{src/cfg/final}

\RemoveFromHook{file/config.tex/after}[aulatex-rsds]

\usepackage{tabularx}
\usepackage{booktabs}
\usepackage{amsmath,amssymb}
\usepackage[backend=biber,style=apa,sorting=nyt]{biblatex}
\definecolor{uanlblue}{RGB}{0,55,108}
\definecolor{tableblue}{RGB}{233,240,247}
\def\portraittitlecolor{uanlblue}
\def\sectioncolor{uanlblue}
\def\ssectioncolor{uanlblue}
\def\documenttitlehf{Responsabilidad Social: Evidencia 3 --- Simulación de Conflicto Socioambiental}
\hypersetup{colorlinks=true,linkcolor=uanlblue,urlcolor=uanlblue,citecolor=uanlblue,pdfauthor={Carlos Mauricio Carvajal Coronado},pdftitle={Simulación de conflicto socioambiental por trasvase de agua},pdfsubject={Evidencia 3 - RSyDS UANL}}
\addbibresource{UANL/ingeniero-agronomo/responsabilidad-social-y-desarrollo-sustentable/responsabilidad-social-y-desarrollo-sustentable.bib}
\newcolumntype{Y}{>{\raggedright\arraybackslash}X}
\setlength{\emergencystretch}{2.5em}

\begin{document}
\hypersetup{pageanchor=false}
\begingroup
\fontsize{14}{17}\selectfont
\renewcommand{\Huge}{\fontsize{14}{17}\selectfont\bfseries}
\renewcommand{\large}{\fontsize{14}{17}\selectfont}
\templatePortrait
\clearpage
\endgroup
\hypersetup{pageanchor=true}
\templatePagecfg
\singlespacing
\setlength{\tabcolsep}{8pt}
\renewcommand{\arraystretch}{1.18}

\section*{Resumen}
En el presente trabajo se formaliza la Evidencia 3 de la Unidad de Aprendizaje de Responsabilidad Social y Desarrollo Sustentable mediante un reporte analítico de Simulación y Juego de Roles en torno a uno de los conflictos socioambientales y agrarios más complejos del noreste de México: la disputa por el agua almacenada en la Presa El Cuchillo (China, Nuevo León) y su trasvase al Distrito de Riego 026 (Bajo Río San Juan, Tamaulipas). Se estructuran cuatro roles representativos (productores agrícolas del DR 026, organismo operador metropolitano SADM, autoridad federal CONAGUA y comunidades ejidales defensoras del caudal ecológico). A través de una metodología de mediación deliberativa fundamentada en la norma ISO 26000 y el desarrollo sustentable, se evalúan las tensiones entre seguridad alimentaria y derecho humano al agua, formulando un acuerdo vinculante de modernización hídrica, intercambio de aguas residuales tratadas y compensación ambiental de cuenca.

\vspace{0.4cm}
\noindent\textbf{Palabras clave:} Juego de roles, conflicto socioambiental, trasvase de agua, Distrito de Riego 026, ISO 26000, gobernanza del agua.

\section{Contexto del Conflicto Socioambiental Agrario}
En el marco del Convenio de Coordinación de 1996 entre el Gobierno Federal, el Gobierno de Nuevo León y el Gobierno de Tamaulipas, se estableció que si la Presa El Cuchillo cuenta con más de $315\text{ millones de m}^3$ al 31 de octubre y la Presa Marte R. Gómez en Tamaulipas almacena menos de $700\text{ millones de m}^3$, Nuevo León debe trasvasar el excedente para abastecer a más de 4,000 usuarios hidroagrícolas que cultivan 76,000 hectáreas de granos básicos en el DR 026 \parencite{conagua2023}.
Sin embargo, el crecimiento metropolitano de Monterrey, la sequía recurrente y la construcción del acueducto El Cuchillo II han intensificado las presiones políticas y sociales para restringir el trasvase, provocando severas fricciones intersectoriales y agravando la incertidumbre de los agricultores de la cuenca baja \parencite{cantu2016}.

\section{Diseño Metodológico del Juego de Roles}
La actividad formativa asignó cuatro posturas antagónicas y complementarias a equipos de trabajo:

\subsection{Rol 1: Productores Agrícolas del Distrito de Riego 026 (Tamaulipas)}
\begin{itemize}
    \item \textbf{Postura:} Exigencia de cumplimiento estricto del Convenio de 1996 y salvaguarda de sus derechos de concesión vigentes.
    \item \textbf{Argumentos:} El sector agrícola genera alimentos estratégicos nacionales (maíz y sorgo). No cumplir el trasvase implica quiebra económica regional y vulneración de la certeza jurídica agraria.
\end{itemize}

\subsection{Rol 2: Servicios de Agua y Drenaje de Monterrey (Nuevo León)}
\begin{itemize}
    \item \textbf{Postura:} Prioridad ineludible al abasto de uso público urbano para más de 5.3 millones de ciudadanos.
    \item \textbf{Argumentos:} El Artículo 4to Constitucional y la Ley de Aguas Nacionales establecen la primacía del consumo doméstico ante escenarios de emergencia y estiaje prolongado.
\end{itemize}

\subsection{Rol 3: Comisión Nacional del Agua (CONAGUA - Árbitro Federal)}
\begin{itemize}
    \item \textbf{Postura:} Imparcialidad técnica, balance hídrico de cuenca y aplicación del marco regulatorio federal.
    \item \textbf{Argumentos:} Obligación de preservar la estabilidad hídrica integral, evitar conflictos de orden público y vigilar la capacidad de almacenamiento de los vasos de captación.
\end{itemize}

\subsection{Rol 4: Comunidades Ejidales Ribereñas y Ecologistas de la Cuenca}
\begin{itemize}
    \item \textbf{Postura:} Garantía del caudal ecológico mínimo del Río San Juan y restauración de humedales ribereños.
    \item \textbf{Argumentos:} Las presas y extracciones han secado tramos enteros del río, extinguiendo fauna nativa y degradando las norias ejidales comunitarias.
\end{itemize}

\section{Matriz de Negociación y Acuerdos de Responsabilidad Social}
Mediante una sesión de mesa redonda deliberativa guiada por los principios de la responsabilidad social (rendición de cuentas, transparencia, respeto a los derechos humanos y a la legalidad según la ISO 26000), se consensuó el siguiente acuerdo:

\begin{table}[htbp]
\centering
\small
\caption{Matriz de acuerdos de la simulación del conflicto hídrico El Cuchillo -- DR 026.}
\label{tab:acuerdo-cuchillo}
\begin{tabularx}{\textwidth}{@{}l Y Y@{}}
\toprule
\textbf{Eje de Compromiso} & \textbf{Acción Concertada} & \textbf{Responsable / Seguimiento} \\
\midrule
\textbf{Tecnificación del Riego} & Inversión conjunta federal-estatal de \$1,800 mdp para entubamiento de canales secundarios en el DR 026, reduciendo mermas por infiltración en $35\%$. & CONAGUA / Módulos de Riego DR 026. \\
\addlinespace
\textbf{Retorno de Agua Tratada} & Cumplimiento estricto del compromiso de devolver al río San Juan $189\text{ millones de m}^3\text{/año}$ de agua tratada con calidad NOM-001 desde las plantas Dulces Nombres y Pesquería. & SADM / SEMARNAT / PROFEPA. \\
\addlinespace
\textbf{Caudal Ecológico y Reforestación} & Establecimiento de un gasto ecológico continuo de $2.5\text{ m}^3\text{/s}$ y fondo ambiental para reforestar 15,000 ha de la cuenca alta con pináceas nativas. & Fideicomiso de Cuenca / Comunidades Ejidales / UANL. \\
\bottomrule
\end{tabularx}
\end{table}

\section{Conclusiones}
1. El ejercicio de simulación demostró que los conflictos socioambientales no pueden resolverse mediante imposiciones unilaterales ni judicialización polarizante, sino a través del diálogo informado y la corresponsabilidad social.
2. Como profesionales de la agronomía, es imprescindible integrar la visión hidrológica y de cuenca, comprendiendo que la eficiencia en el uso del agua en parcelas es el verdadero garante de la paz y sustentabilidad regional.

\newpage
\printbibliography[title={Referencias bibliográficas}]
\end{document}

